\documentclass[10pt]{article}

\usepackage[utf8]{inputenc}

\usepackage[T1]{fontenc}

\usepackage{amsmath}

\usepackage{amsfonts}

\usepackage{amssymb}

\usepackage[version=4]{mhchem}

\usepackage{stmaryrd}

\usepackage{bbold}



\begin{document}

СУЩЕСТВЕННОЕ ОТОБРАЖЕНИЕ - непрерЫвное отображение $f$ топологич. пространства $X$ в открытый симплекс $T^{n}$ такое, что всякое непрерывное отображение $f_{1}: X \rightarrow T^{n}$, совпадающее с $f$ во всех точках множества $f^{-1}\left(\overline{T^{n}} \backslash T^{n}\right)$, есть отображение на все $T^{n}$. Напр., тождественное отображение $\overline{T^{n}}$ на себя есть С. о.



Jum.: [1] А л е к с а н д р о в П. С., П а с ын к о в Б. А., Введение в теорию размерности ..., М., 1973.



М. И. Войцеховский. СУЩЕСТВОВАНИЯ КВАНТОР - логическая операция, служащая для образования высказываний с помощью оборота «для некоторых $x$ » («существует $x$, для которого》). В формализованных языках С. к. обозначается через $\exists x_{r}(\exists x), \mathrm{U}_{x}, \vee_{x}, \Sigma_{x}$ В. Е. Плиско.



СФЕР ГОМОТОПИЧЕСКИЕ ГРУППЫ - объекТ ИзУчения классич. теории гомотопий. Вычисление С. г. г. $\pi_{i}\left(S^{n}\right)$ в свое время (особенно в $50-\mathbf{x}$ гг.) рассматривалось как одна из центральных задач топологии. Топологи надеялись, тто әти грушпы удастся полностью вычислить и что с их помощью можно будет решать другие классификационные гомотопич. задачи. Эти надежды в основном не сбылись: С. Г. г. удалось вычислить лишь частично, и с развитием теории обобщенных когомологий задача их вычисления стала менее актуальной. Все же накопленная информация об этих группах не пропала даром, она нашла применения там, где их не ждали, в частности в дифференциальной топологии (классификации дифференциальных структур на сферах и многомерных узлов).



I. О бщ ая т е о рия. 1) Если $i<n$ или $i>n=1$, To $\pi_{i}\left(S^{n}\right)=0$.



\begin{enumerate}

  \setcounter{enumi}{1}

  \item $\pi_{n}\left(S^{n}\right)=\mathbb{Z} \quad($ т е о р м а $\quad$ Б ф а); әтот изоморфизм относит элементу группы $\pi_{n}\left(S^{n}\right)$ степень представляющего его отображения $S^{n} \rightarrow S^{n}$.



  \item Группы $\pi_{4 m-1}\left(S^{2 m}\right)$ имеют ранг 1; прочие группы $\pi_{i}\left(S^{n}\right)$ с $i \neq n$ конечны.



\end{enumerate}



Гомоморфизм надстройки



\[

E: \pi_{i}\left(S^{n}\right) \longrightarrow \pi_{i+1}\left(S^{n+1}\right)

\]



относит әлементу группы $\pi_{i}\left(S^{n}\right)$, представляемому сфероидом $f: S^{i} \rightarrow S^{n}$, класс сфероида $E f: S^{i+1} \rightarrow S^{n+1}$, определяемого формулой



\[

E f\left(\sqrt{1-x^{2}} x, x\right)=\left\{\begin{array}{l}

\left(\sqrt{1-x^{2}} f(x), x\right),|x| \leqslant 1, \\

(0, x),|x|=1,

\end{array}\right.

\]



где $x \in S^{i}, x \in \mathbb{R}$



\begin{enumerate}

  \setcounter{enumi}{3}

  \item Гомоморфизм $E$ является изоморфизмом при $i>2 n-$ -1 и әпиморфизмом при $i \geqslant 2 n-1$.

\end{enumerate}



Таким образом, при каждом $k$ группы $\pi_{n+k}\left(S^{n}\right)$ могут быть составлены в последовательность



\[

\pi_{1+k}\left(S^{1}\right) \stackrel{E}{\longrightarrow} \pi_{2+k}\left(S^{2}\right) \stackrel{E}{\longrightarrow} \pi_{3+k}\left(S^{3}\right) \stackrel{E}{\longrightarrow \ldots,}

\]



в $(k+2)$-м члене к-рой наступает стабилизация; группа $\pi_{n+k}\left(S^{n}\right)$ с $n \geqslant k+2$ наз. $k$-й с т а б и л ь н о й С. г. г. и обозначается $\pi_{k}^{\mathrm{s}}$. При этом $\pi_{k}^{s}=0$ при $k<0$ и $\pi_{0}^{s}=\mathbb{Z}$.



Как и в гомотопических гpynпах любого топологич. пространства, в С. г. г. определено умножение Уайтхеда:



$\pi_{i}\left(S^{n}\right) \times \pi_{j}\left(S^{n}\right) \longrightarrow \pi_{i+j-1}\left(S^{n}\right),(\alpha, \beta) \longrightarrow[\alpha, \beta]$.



$\kappa$ обычным его свойствам (дистрибутивность, косая коммутативность, тождество Якоби) добавляется



\begin{enumerate}

  \setcounter{enumi}{4}

  \item $E[\alpha, \beta]=0$.

\end{enumerate}



Умножение Уайтхеда позволяет сделать следующее уточнение к 4):



\begin{enumerate}

  \setcounter{enumi}{5}

  \item ядро эпиморфизма $E: \pi_{2 n-1}\left(S^{n}\right) \rightarrow \pi_{2 n}\left(S^{n+1}\right)$ порождается классом $\left[i_{n}, i_{n}\right]$, где $i_{n}-$ каноническая образующая группы $\pi_{n}\left(S^{n}\right)$ (представляемая тождественным сфероидом).

\end{enumerate}



С умножением Уайтхеда тесно связан Хопба инвариант $H(\alpha)$, определенный для $\alpha \in \pi_{4 m-1}\left(S^{2 m}\right)$. Так, элемевт группы $\pi_{3}\left(S^{2}\right)$, представляемый отображением Хопфа $h: S^{3} \rightarrow S^{2}$, действующим по формуле $h\left(z_{1}, z_{2}\right)=$ $=z_{1}: z_{2}$ (в к-рой $S^{3}$ интерпретируется как единичная сфера пространства $\mathbb{C}^{2}$, а $S^{2}-$ как $\mathbb{C} P^{1}$ ), имеет инвариант Хопфа, равный 1



\begin{enumerate}

  \setcounter{enumi}{6}

  \item Отображение $H: \pi_{3}\left(S^{2}\right) \rightarrow \mathbb{Z}$ есть изоморфизм.



  \item $H\left(\left[i_{2 m}, i_{2 m}\right]\right)=2$.



\end{enumerate}



Следствием 8) является бесконечность групп $\pi_{4 m-1}\left(S^{2 m}\right)$, уже утверждавшаяся в 3$)$.



\begin{enumerate}

  \setcounter{enumi}{8}

  \item При $m \neq 1,2,4$ в $\pi_{4 m-1}\left(S^{2 m}\right)$ отсутствуют элементы с нететным инвариантом Хопфа (как было известно задолго до доказательства этой теоремы, ее утверждение равносильно следующей г йо те зе Ф р об е н и у с а: при $l \neq 1,2,4,8$ в $\mathbb{R}^{l}$ отсутствует билинейное умножение с однозначным делением на ненулевые әлементы).

\end{enumerate}



Специфическим для сфер является к о м по з и ц ион но е умножен и е $\pi_{i}\left(S^{j}\right) \times \pi_{j}\left(S^{n}\right) \longrightarrow \pi_{i}\left(S^{n}\right),(\beta, \alpha) \longrightarrow \alpha \circ \beta$,



определяемое при помощи компонирования представляющих отображений.



\begin{enumerate}

  \setcounter{enumi}{9}

  \item Для любых $\alpha, \alpha_{1}, \alpha_{2} \in \pi_{j}\left(S^{n}\right), \beta, \beta_{1}, \beta_{2} \in \pi_{i}\left(S^{j}\right)$, $\delta \in \pi_{i-1}\left(S^{j-1}\right), \gamma \in \pi_{k}\left(S^{j}\right)$, имеет место:

\end{enumerate}



a) $(\alpha \circ \beta) \circ \gamma=\alpha \circ(\beta \circ \gamma)$,



б) $\alpha \circ\left(\beta_{1}+\beta_{2}\right)=\alpha \circ \beta_{1}+\alpha \circ \beta_{2}$,



в) $\left(\alpha_{1}+\alpha_{2}\right) \circ E \delta=\alpha_{1} \circ E \delta+\alpha_{2} \circ E \delta$



г) $E(\alpha \circ \beta)=E \alpha \circ E \beta$.



«Левый закон дистрибутивности» $\left(\alpha_{1}+\alpha_{2}\right) \circ \beta=\alpha_{1} \circ \beta+$ $+\alpha_{2} \circ \beta$, вообще говоря, места не имеет. Утверждение г) позволяет определить с т а б и ль н о е ко мпозй ци онно е умножени е



\[

\pi_{q}^{s} \times \pi_{r}^{s} \rightarrow \pi_{q+r}^{s},(\beta, \alpha) \longrightarrow \alpha \circ \beta .

\]



\begin{enumerate}

  \setcounter{enumi}{10}

  \item Для любых $\alpha, \alpha_{1}, \alpha_{2} \in \pi_{r}^{s}, \beta, \beta_{1} \beta_{2} \in \pi_{q}^{s}, \gamma \in \pi_{p}^{s}$ имеют место утверждения а), б) из (0), а также:

\end{enumerate}



B $\left.^{\prime}\right)\left(\alpha_{1}+\alpha_{2}\right) \circ \beta=\alpha_{1} \circ \beta+\alpha_{2} \circ \beta$,



д) $\alpha \circ \beta=(-1)^{q r} \beta \circ \alpha$.



II. М е т о ды вы ч и с л ен и я. Геометрич. метод Понтрягина (см. [1]), предложенный в сер. 30-х гг., основывается на следующем определении. Гладкое $m$-мерное компактное многообразие $X$ пространства $\mathbb{R} i$ наз. о с в а щ е н н ы м, если ва нем задано гладкое поле трансверсальных к нему $(i-m)$-реперов; само это поле ваз. о с в а щ ө в и е м. Оснащенные многообразия $X_{0}, X_{1} \subset \mathbb{R} i$, нө имеющие края, наз. к о б о рд а н т н ы м и, если существует оснащенное многообразие $Y \subset \mathbb{R} i \times(0,1] \subset \mathbb{R}^{i+1} \mathrm{c} \partial Y=\left(X_{0} \times 0\right) \cup\left(X_{1} \times 1\right)$, у к-рого сужение оснащения на $X_{0} \times 0$ и $X_{1} \times 1$ содержится в $\mathbb{R} i \times 0$ и $R i \times 1$, и при естественном отождествлении $\mathbb{R} i \times 0$ и $\mathbb{R} i \times 1$ с. $\mathbb{R} i$ превращается в заданное оснащение многообразий $X_{0}$ и $X_{1}$. Множество классов кобордантных оснащенных $m$-мерных многообразий без края в $R^{i}$ обозначается через $\Omega^{m}(i)$.



\begin{enumerate}

  \item Имеет место взаимно однозначное соответствие между $\pi_{i}\left(S^{n}\right)$ и $\Omega^{i-n}(i)$.

\end{enumerate}



Метод дает хорошие результаты при небольших $i-n$. Он позволяет также доказать нек-рые из теорем п. I и доставляет разнообразную геометрич. информацию о многообразиях малых размерностей.



Другую группу методов составляют элементарные алгебраич. методы, состоящие в использовании гомотопич. последовательностей различных расслоений, свойств умножения Уайтхеда, композиционного умножения и соответствующего высшего умножения (скобок Тоды, см. [3]), а также следующую т е о р е м у Д же й м с а.



\begin{enumerate}

  \setcounter{enumi}{1}

  \item Имеется последовательность групп и гомоморфизмов

\end{enumerate}



\[

\begin{aligned}

\ldots \pi_{i}\left(S^{n}\right) & \stackrel{E}{\longrightarrow} \pi_{i+1}\left(S^{n+1}\right) \stackrel{H}{\longrightarrow} \pi_{i+1}\left(S^{2 n+1}\right) \stackrel{P}{\longrightarrow} \\

& \longrightarrow \pi_{i-1}\left(S^{n}\right) \longrightarrow \ldots,

\end{aligned}

\]





\end{document}